\section{The determinant construction}
\label{s2-construction}

This section follows Section~2 of the source, pp.~3--5.  We change the
functional to one adapted to $\zeta(9)$, while retaining the source's
matrix size, numerator $D_N^6$, and scalar normalization.  The degree and
positive integral representation can then be calculated directly.  The
arithmetically normalized decay statement is a separate conditional target.

\subsection{The functional and the matrix}
\label{s2-functional}

As in source (2.1), put
\begin{equation}
 K=40n,\qquad N=3n,\qquad h=K-N=37n,\qquad
 \alpha=\frac3{40},\quad\lambda=\frac{37}{40},\quad
 H=1+2\alpha=\frac{23}{20}.
 \label{s2-parameters}
\end{equation}
Define $D_0(t)=1$ and
\[
 D_m(t)=\prod_{j=1}^{m}(t+j^2),\qquad
 H_j^{(9)}=\sum_{v=1}^{j}\frac1{v^9},\qquad H_0^{(9)}=0.
\]
Use Bernoulli numbers with $B_1=-1/2$.  The analogues of source
(2.2)--(2.3) are
\begin{align}
 \mu_9(t^e)
 &=\frac{(-1)^e B_{2e+2}}{40320}
       \prod_{r=3}^{9}(2e+r),\qquad e\geq0,
 \label{s2-monomials}\\
 \mu_{9,X}\!\left(\frac1{t+j^2}\right)
 &=j^8\bigl(X-H_j^{(9)}\bigr)-\frac18+\frac1{2j},
       \qquad j\geq1.
 \label{s2-poles}
\end{align}
Here $40320=8!$, and $\mu_{9,X}$ agrees with $\mu_9$ on polynomials.
For example, $\mu_9(1)=3/4$ and $\mu_9(t)=11/8$; these provide simple
checks on the factorial.  Polynomial division and simple partial
fractions uniquely extend the functional to rational functions whose
poles are simple and lie among the negative integer squares.  Only the
pole contribution depends on $X$, and it does so affinely.
For instance, the first two pole values are
\[
 \mu_{9,X}\!\left(\frac1{t+1}\right)=X-\frac58,
 \qquad
 \mu_{9,X}\!\left(\frac1{t+4}\right)=256X-\frac{2051}{8}.
\]
The factor $256=2^8$ and the harmonic increment $2^{-9}$ must be changed
together.  Evaluating the old fifth-order functional at $X=\zeta(9)$
would retain its old moments and would not yield the weight constructed
below.  Both parts of the functional, not just its argument, are being
replaced.

Retaining source (2.4)--(2.5), define
\begin{equation}
 G_{9,K}(X)=
 \left[\mu_{9,X}\!\left(\frac{D_N(t)^6t^{i+j}}{D_K(t)}\right)
 \right]_{0\leq i,j<h},\qquad
 \Delta_{9,K}(X)=\det G_{9,K}(X),
 \label{s2-matrix}
\end{equation}
and
\begin{equation}
 S_K=\frac{(K!)^{2h}4^{h-1}}
 {(N!)^{12h}\displaystyle\prod_{i=1}^{h-1}((2i)!)^2},
 \qquad F_{9,K}(X)=S_K\Delta_{9,K}(X).
 \label{s2-normalization}
\end{equation}
The numerator remains $D_N^6$: the order of the target zeta value alone
does not require a different numerator power.  This choice also preserves
the leading real external field used later.  For fixed $M\geq40$, the
analogue of source (2.6) is
\begin{equation}
 Q_{9,K,M}(X)=m_{9,K,M}F_{9,K}(X),
 \qquad m_{9,K,M}\in\mathbb Q_{>0},
 \label{s2-Q}
\end{equation}
where Section~5 must supply the prime exponents and their cost.
Cancellation of $D_N$ leaves a denominator with distinct linear factors,
so all matrix entries lie in the stated domain.  Their polynomial
quotients can have substantial degree; they are retained in the
definition and cannot be discarded during the later denominator
calculation.  The factor $S_K$ is positive and independent of $X$.
Its reuse fixes the normalization being compared across the two zeta
orders, but by itself makes no assertion that every coefficient is
integral.

\subsection{The polynomial estimates: conditional targets}
\label{s2-targets}

\noindent\textbf{Draft analogue of Theorem 2.1 [Conditional].}
For fixed $M\geq40$ and $K\in40\mathbb Z_{>0}$ with $K\geq200M^2$,
the tasks for the later sections are to establish a positive rational
normalizer in \eqref{s2-Q} for which
\begin{equation}
 Q_{9,K,M}\in\mathbb Z[X],\qquad
 \deg Q_{9,K,M}=h,\qquad Q_{9,K,M}(\zeta(9))>0,
 \label{s2-integer-target}
\end{equation}
and to control its value by comparing the estimates
\begin{align}
 \log m_{9,K,M}&\leq A_{9,M}K^2+o_M(K^2),
 \label{s2-arithmetic-target}\\
 \log F_{9,K}(\zeta(9))&\leq U K^2+O(K\log K),
 \qquad U=-\frac{2733991}{2000000}.
 \label{s2-real-target}
\end{align}
If these estimates hold, their sum gives
\begin{equation}
 \log Q_{9,K,M}(\zeta(9))
 \leq (A_{9,M}+U)K^2+o_M(K^2).
 \label{s2-combined-target}
\end{equation}
A quadratic decay conclusion additionally requires $A_{9,M}+U<0$.
This sign is a goal, not an assumption or a conclusion of this section.
The numerical bounds in source (2.7)--(2.8) therefore have no asserted
$\zeta(9)$ counterparts here.  Degree and positivity will follow below
for every positive normalizer; coefficient integrality and the sign of
the combined estimate require the later arithmetic calculation.
In particular, a larger denominator-clearing multiplier can preserve
all three properties in \eqref{s2-integer-target} while destroying the
small-value estimate.  This is why the draft records its quadratic cost
as a separate quantity.  A favorable real integral estimate and exact
nonvanishing do not determine that cost.  Throughout this comparison,
$M$ is fixed before $K$ tends to infinity. The prime-sum argument in
Section~5 supplies a limiting quadratic bound, not an explicit
$O_M(K\log K)$ remainder for the arithmetic part.

\subsection{Degree}
\label{s2-degree}

\noindent\textbf{[Derived].}
Let $D_{\mathrm{tail}}=D_K/D_N$ and $W=D_N^5$.  Cancellation in a
matrix entry gives $Wt^{i+j}/D_{\mathrm{tail}}$.  At the pole $-k^2$,
its residue is
\[
 \frac{W(-k^2)(-k^2)^{i+j}}{D_{\mathrm{tail}}'(-k^2)}.
\]
Consequently, for the square Vandermonde matrix
$V_{ij}=(-j^2)^i$, indexed by $0\leq i<h$ and $N<j\leq K$,
\begin{equation}
 [X]G_{9,K}(X)=V\operatorname{diag}\!\left(
 \frac{j^8D_N(-j^2)^5}{D_{\mathrm{tail}}'(-j^2)}
 \right)_{N<j\leq K}V^{\mathsf T}.
 \label{s2-leading-matrix}
\end{equation}
There are exactly $h$ surviving pole locations.  If $x_j=-j^2$, then
\[
 \prod_jD_{\mathrm{tail}}'(x_j)
   =(-1)^{h(h-1)/2}\prod_{a<b}(x_b-x_a)^2.
\]
The Vandermonde square thus cancels the denominator in the determinant
of \eqref{s2-leading-matrix}.  The analogue of source (2.9) is
\begin{equation}
 [X^h]\Delta_{9,K}(X)
   =(-1)^{h(h-1)/2}
     \prod_{j=N+1}^{K}j^8D_N(-j^2)^5\ne0.
 \label{s2-leading-coefficient}
\end{equation}
Each entry is affine, so the determinant has degree at most $h$; this
nonzero coefficient makes its degree exactly $h$.  Multiplication by
$S_K$ or a positive rational $m_{9,K,M}$ preserves that degree.

\subsection{The positive integral representation}
\label{s2-positivity}

\noindent\textbf{Analogue of Proposition 2.2 [Derived].}
For every rational function in the domain of $\mu_{9,X}$,
\begin{equation}
 \mu_{9,\zeta(9)}(R)=\int_0^\infty R(y^2)w_9(y)\,dy,
 \qquad
 w_9(y)=\frac{(2\pi)^8y^9}{20160}
           \sum_{\ell\geq1}\ell^8e^{-2\pi\ell y}>0.
 \label{s2-integral}
\end{equation}
This specializes the general odd-order weight
$2(2\pi)^{s-1}y^s\sum_{\ell\geq1}\ell^{s-1}e^{-2\pi\ell y}/(s-1)!$
at $s=9$.  The specialization is useful because it changes the order
without introducing lower odd zeta values.

For monomials, integrating the nonnegative series term by term gives
\begin{align}
 \int_0^\infty y^{2e}w_9(y)\,dy
 &=\frac{(2e+9)!}{20160(2\pi)^{2e+2}}\zeta(2e+2)
 \notag\\
 &=\frac{(-1)^eB_{2e+2}}{40320}
       \frac{(2e+9)!}{(2e+2)!}
 =\mu_9(t^e).
 \label{s2-moment-integral}
\end{align}
The second equality uses Euler's even-zeta formula, as in source
Proposition~2.2.

For the pole calculation set $f(y)=(e^{2\pi y}-1)^{-1}$.  Since
$w_9(y)=y^9f^{(8)}(y)/20160$, the required derivative is obtained by
first dividing:
\[
 \frac{y^9}{y^2+a^2}
 =y^7-a^2y^5+a^4y^3-a^6y
    +a^8\frac{y}{y^2+a^2}.
\]
The polynomial part disappears after eight derivatives, while
$y/(y^2+a^2)=\operatorname{Im}(a-iy)^{-1}$.  Therefore
\begin{equation}
 \left(\frac{d}{dy}\right)^8\frac{y^9}{y^2+a^2}
   =8!\,a^8\operatorname{Im}(a-iy)^{-9}.
 \label{s2-eighth-derivative}
\end{equation}
Eight integrations by parts now give
\begin{equation}
 \int_0^\infty\frac{w_9(y)}{y^2+a^2}\,dy
   =2a^8\int_0^\infty
       \frac{\operatorname{Im}(a-iy)^{-9}}{e^{2\pi y}-1}\,dy.
 \label{s2-pole-ibp}
\end{equation}
To check the boundary terms, write $g(y)=y^9/(y^2+a^2)$.
For $0\leq r\leq7$, the product $g^{(7-r)}f^{(r)}$ is
$O(y^{r+2})O(y^{-r-1})=O(y)$ at zero.  At infinity the exponential
decay of every derivative of $f$ absorbs the polynomial factors.

Hermite's formula, the identity cited by the source as DLMF
(25.11.29), reads in this case
\[
 \zeta(9,a)=\frac1{2a^9}+\frac1{8a^8}
   +2\int_0^\infty
       \frac{\operatorname{Im}(a-iy)^{-9}}{e^{2\pi y}-1}\,dy,
 \qquad a>0.
\]
Thus \eqref{s2-pole-ibp} becomes
\begin{equation}
 \int_0^\infty\frac{w_9(y)}{y^2+a^2}\,dy
   =a^8\zeta(9,a)-\frac1{2a}-\frac18.
 \label{s2-hurwitz-pole}
\end{equation}
At $a=j\in\mathbb Z_{>0}$, substitute
$\zeta(9,j)=\zeta(9)-H_j^{(9)}+j^{-9}$ to recover exactly
\eqref{s2-poles}.  Together with the monomial calculation, linearity
proves \eqref{s2-integral}.  Moreover $w_9(y)\to1/\pi$ at zero and
$w_9(y)=O((1+y)^9e^{-2\pi y})$ at infinity, so every integral used
here converges.
The endpoint statements can also be read directly from the weight.
Near zero the Laurent expansion of $f$ gives
$f^{(8)}(y)=8!/(2\pi y^9)+O(1)$, which cancels the factor $y^9$.
For $y\geq1$,
\[
 \sum_{\ell\geq1}\ell^8e^{-2\pi\ell y}
 \leq e^{-2\pi y}
      \sum_{\ell\geq1}\ell^8e^{-2\pi(\ell-1)},
\]
and the final sum is a finite constant.  Thus rational functions in the
domain are bounded near the integration endpoint zero and have at most
polynomial growth at infinity, which the exponential decay absorbs.

Finally, if $q$ is a nonzero real polynomial of degree below $h$, then
\begin{equation}
 \int_0^\infty
   \frac{D_N(y^2)^6}{D_K(y^2)}q(y^2)^2w_9(y)\,dy>0.
 \label{s2-positive-quadratic-form}
\end{equation}
This is the quadratic form of $G_{9,K}(\zeta(9))$ at the coefficient
vector of $q$.  All weight factors are positive for $y>0$, and $q(y^2)$
has only finitely many zeros.  Hence the matrix is positive definite,
$\Delta_{9,K}(\zeta(9))>0$, and every positive normalization retains
nonvanishing.  This completes the local algebraic construction without
assuming the later coefficient-clearing cost is small enough.
